Requirements of Approach Section (2000 words maximum)

Explain how you have designed your work so that it:
is effective and appropriate to achieve your objectives
is feasible, and comprehensively identifies any risks to delivery and how you will manage them
uses a clearly written and transparent methodology (if applicable)
summarises the previous work and describes how you will build on and progress this work (if applicable)
will maximise translation of outputs into outcomes and impacts
Within the Approach section we also expect you to:
demonstrate access to the appropriate services, facilities, infrastructure, or equipment to deliver the proposed work
provide a detailed and comprehensive project plan, including milestones and timelines in the form of a Gantt chart or diagram
explain who you intend to collaborate with at the host organisation and your plans for wider research collaborations (project partner details should be provided only in that section)
include details of work that will take place as part of the proposed fellowship at a second UK or overseas organisation (if applicable)
explain and justify how you will approach diversity and inclusion in the study population and follow the MRC embedding diversity in research design policy (if applicable)
show how you will use male and female animals or tissues and cells from female and male donors (if applicable) in your research. If you are not proposing to do this, justify why
explain and justify the inclusion of public partnerships (if applicable) and the added value these offence

Including ‘References’ and ‘Images’
You are permitted to include ‘References’ and ‘Images’ within this section. If you choose to include these, you must refer to the guidance within the ‘Vision’ section. Failure to follow the guidance may result in your application being rejected.
You should detail any information regarding your statistical design in the 'Reproducibility and statistical design' section.


_______________________________


3 Approach

3.1 Strategic Overview: From Mechanistic Evolution to Personalised Prediction

This Professional Fellowship establishes a comprehensive workflow for predictive pathogen genomics, systematically moving from fundamental mechanistic understanding to the refinement of advanced machine learning models, and culminating in clinical validation. Building on the lab's established foundation in chronic pathogen genomics, the project addresses the limitations that currently prevent the translation of these discoveries into a universally deployable clinical tool. The proposed workflow is structured to systematically bridge this gap, ensuring we generate deep mechanistic insights, a testable pipeline of benchmarks, and a model ready for clinical testing and translation to patient care.
Theme 1: Foundational Mechanistic Genomics. A comprehensive analysis into the unique pathoadaptive mechanisms of the Enterobacteriaceae Klebsiella pneumoniae, focusing on rampant horizontal gene transfer (HGT) and structural variation—mechanisms distinct from those previously studied by the group. This work identifies critical factors that need to be incorporated into models in themes 2 and 3.
Theme 2: Machine Learning Refinement with Bacformer. A rigorous evaluation of the state-of-the-art Bacformer model against the new mechanistic insights from Theme 1, explicitly identifying and addressing its limitations related to HGT, plasmids, and non-coding DNA to develop a next-generation hybrid model.
Theme 3: Clinical Validation and Hybrid Outcome Model. A critical translational step that uses the refined genomic prediction score from Theme 2, combines it with rich patient history, and validates its predictive power in the real-world clinical data of the SputaGen prospective study.
3.2 Theme 1: Foundational Mechanistic Genomics of Klebsiella Pathoadaptation

Our group has proven that saltatory changes, new genes, and alterations in transcriptional regulation dominate evolution and pathoadaptation in M. abscessus (Bryant et al., 2021; Bryant et al., 2016) and P. aeruginosa (Weimann et al., ref.), two important chronic infecting agents that are notoriously difficult to treat. We still have a limited understanding of the pathoadaptation of other significant threats to global lung health. Klebsiella species (and K. pneumoniae in particular) are one of the ESCAPPM pathogens and are a significant cause of mortality worldwide, with rapidly escalating rates of AMR (reference needed). Like E. coli and A. baumannii, it is critical we understand how this Enterobacterale evolves to escape immunity, become invasive (virulent), and avoid antibiotics.

Klebsiella represents a fundamentally different challenge from the Mycobacterium and Pseudomonas species we have primarily studied. It is characterised by rampant Horizontal Gene Transfer (HGT) and genetic recombination events, notably in the capsule (Holt, 2015 reference) and the acquisition of large plasmids (Cazares et al., ref.) carrying both AMR and virulence genes. While the mechanisms driving the global success of subtypes like Sequence Lineage (SL) 258 (which accounts for over 20% of global samples) remain unclear, they are hypothesised to involve a complex interplay of HGT, plasmids, and rapid phenotypic switching via insertion transposable elements, mechanisms quite distinct from the transcriptional hotspot changes seen in M. abscessus and P. aeruginosa.

Data Foundation

The first part of this fellowship will conduct a foundational genomic study on Klebsiella. This work will be grounded in a rigorously curated, large-scale dataset:
Global Klebsiella Genome Collection: A collection of approximately 80,000 high-quality, whole-genome assemblies of the Klebsiella pneumoniae species complex (KpSC). This collection is globally distributed, with significant sampling in Africa and Asia, and has extensive metadata curation, providing isolation source, host annotations, collection dates, and countries for ~80% of samples.
Deep Phenotypic Data (EBI AST): We will integrate a large-scale antibiotic susceptibility testing (AST) dataset comprising 7,000 Klebsiella isolates with over 140,000 AST results. Critically, 23 antibiotics have sufficient data (>~2,500 AST results each) to represent the full spectrum of clinically important treatments used internationally.
Mechanistic Analysis

My bioinformatics group, including Wiemann, Bruchmann, and Dinan, will lead the analysis of variant calling and hotspot changes in transcription, as well as the dating of the emergence of virulent clones. My specific focus will be leading the analysis of plasmids and horizontal gene transfer (HGT), including insertion sequences (ISEscan, MGEfinder), which may be a fundamentally different method of pathoadaptation and immune/antibiotic evasion compared to M. abscessus and P. aeruginosa. This foundational work will allow us to understand the breadth of pathogen evolution and behavior in one of the most important Enterobacterales.

3.3 Theme 2: Machine Learning Refinement with Bacformer

We have already demonstrated the ability of Bacformer to achieve AUROC / AUPRC of > 0.97 on multiple antibiotics reading only from a string of proteins and with no prior knowledge of protein function (Wiatrak, Abelson et al., in prep.). It can also detect phylogeny responsible for resistance and particular genes of resistance using explainability techniques (Captum attribution).

However, in its first iteration (current) it only models protein-coding regions. It is unclear how it handles HGT, plasmids, insertion sequences, and non-coding regulatory regions, all of which are critical to the more complex phenotypes identified in Theme 1. Our group led by Wiatrak is currently iterating Bacformer to create a hybrid protein and DNA large language model that will have representations for both sections, trained on more than 8 billion proteins and over 1.5 million bacterial assemblies.

Part 2 of this fellowship will rigorously test the new model's performance against the mechanistic findings from Theme 1 and refine its structure:
Complex Phenotype Testing: We will thoroughly test Bacformer’s ability to predict complex phenotypes (AMR, Tolerance, and Virulence Potential) in the setting of HGT, including:
Invasive vs. Carriage Prediction: Testing the model's ability to differentiate between strains isolated from invasive disease (Blood/Respiratory/Wound) versus carriage (Stool/Urine) on the $\sim$26,000 curated isolates.
HGT Representation: Examining the model's representations of plasmids and insertion sequences identified by standard bioinformatics in Theme 1. We will test if the model captures plasmids as common embeddings across varied sequence lineages and species, allowing for a better understanding and representation of HGT.
Gene Calling Impact: Analysis of how insertions (identified by ISEscan) impact traditional gene-calling pipelines (Panaroo) and how this affects Bacformer’s performance, leading to the prioritisation of structural improvements in the model's input syntax.
The findings from this theme will directly inform the development of the next-generation Bacformer model, ensuring its foundation is robust enough to handle the genomic complexity found in real-world clinical scenarios.

3.4 Theme 3: Clinical Validation and Translational Impact with SputaGen

This theme provides the critical clinical payoff, leveraging the applicant’s unique role as Principal Investigator (PI) of the SputaGen study to test the predictive tools in a real-world setting.

Bridging the Genomic and Clinical Gap

The core challenge in moving genomic models to the clinic is that treatment outcome is not solely determined by the invading pathogen's genomic potential (its intrinsic resistance or virulence, as predicted by Bacformer). It is profoundly influenced by the patient context: their prior antibiotic exposure, clinical history, and the composition of their persistent microbial reservoir. This is critical because antibiotic-induced treatment failure is often driven by strain replacement from the patient’s own colonizing microbiota, a non-genomic event (Stracy et al., 2022). As demonstrated in UTI, an ML model integrating the patient's historical culture data and clinical factors (e.g., age, catheter use) is a superior predictor of resistance emergence than genotypic data alone.

Therefore, this fellowship is strategically structured: Themes 1 and 2 are essential for refining the genomic component (Bacformer) to accurately predict the organism's inherent behavior (resistance, tolerance, virulence). Theme 3 then serves as the indispensable validation and translational step, requiring the creation of a hybrid clinical outcome model that combines the advanced genomic prediction score from Bacformer with the rich clinical, demographic, and historical data captured prospectively by SputaGen. This hybrid approach is the only way to realize a tool for truly personalised, effective antibiotic prescribing.

Sequencing as the Final Translational Lever

The SputaGen study, designed and led by the applicant, is the ideal validation cohort. The clinical data collection is nearing completion and is fully funded. The essential missing bridge is the pathogen genomic data. The funding requested within this Fellowship for sequencing the SputaGen isolates and/or performing metagenomics is a high-leverage investment ($\sim$£100,000) that immediately unlocks the translational objective of the entire programme. It transforms an existing, mature clinical resource into the final, critical dataset required to build and validate the next-generation clinical prediction tool.

SputaGen Study Design and Resources
Study Status: The SputaGen study (Protocol P03129, ISRCTN 14487574) is a prospective, observational single-centre study at the Cambridge Centre for Lung Infection (CCLI). Sample collection is nearing completion and is fully funded by the Cystic Fibrosis (CF) Foundation Trust, providing a mature and unique validation cohort.
Scale and Scope: The study is designed to capture between 600 and 3,000 IV antibiotic treatment courses over five years. As patients can have multiple treatment courses, this will yield approximately 1,200 to 6,000 antibiotic-patient outcomes for analysis, alongside detailed clinical metadata, conventional microbiology, and Drug Susceptibility Testing (DST). The final analysis will track patient outcomes, including clinical response, patient-reported outcome measures (e.g., CAAT, mMRC), and time to the next acute exacerbation (sAPE).
Sequencing Plan: The Fellowship will include an application for additional sequencing funds to process the SputaGen biobank. This will involve the sequencing of cultured isolates and/or metagenomic sequencing of the sputum samples. This work will be carried out in close collaboration with Dr. Josie Bryant at the Sanger Institute, who provides world-leading expertise in high-throughput sequencing and metagenomic analysis.
Translational Objective

The final phase will be to directly relate the refined Bacformer predictions for AMR, antibiotic tolerance, and virulence potential from Themes 1 and 2 to the SputaGen clinical dataset.
Clinical Outcome Modeling: We will develop clinical outcome models to test whether a combined genomic prediction score (integrating resistance, tolerance, and virulence) derived from the patient's infecting pathogen is a superior predictor of clinical response and treatment failure following intravenous antibiotic courses, independent of conventional AST results.
Tool Translation: The project will culminate in the open-source release of the final data, test pipelines, and models into a reusable “Bacformer-SputaGen toolkit,” establishing a clear pathway for commercial and NHS clinical network engagement.
-----References (Full Reference at Bottom of Document)
(Bryant et al., 2021)
(Bryant et al., 2016)
(Weimann et al., ref.)
(Cazares et al., ref.)
(Boeck, 2023)
(Wiatrak, Abelson et al., in prep.)
(Stracy et al., 2022)
Final response to user: I have completely restructured the Approach section to reflect your new, three-part workflow.
Theme 1 is now a dedicated, foundational mechanistic study on Klebsiella's unique pathoadaptation (HGT, plasmids, insertion sequences), incorporating the large dataset details.
Theme 2 is now a dedicated phase for the refinement and rigorous testing of Bacformer against the mechanistic truths found in Theme 1.
Theme 3 (SputaGen) remains largely the same, clearly positioned as the final, critical clinical validation step requiring the sequencing investment.
Please review the new narrative. What would you like to edit or refine next?


—--------------



Services and Infrastructure / Equipment:
-  High Performance Cluster at University of Cambridge
-  Coding assistant with Cursor AI - I have this on HPC as a part of my workflow
-  Bacformer Hugging Face Models
	-  Klebsiella dataset I have curated
	-  SputaGen dataset
	-  SputaGen sequencing


GANTT CHART WITH TIMELINES

Collaborators:  Josie Bryant at Sanger for Sequencing 
My responsibilities and other organisations:
Diversity, inclusion, ethics and population health:
Animals and gender:  Not relevant
Risks and Limitations: 
Public partnerships: None




